\scanpage{226}
\mspage{142}
% archive sheet number: 225
% element: text | top, continuation and cardinal notation | un cardinal infini
un cardinal infini). Ici, pour la catégorie \deleted{\unsure{la|$\mathcal A$}} $A$, je pose $\operatorname{card}A=\operatorname{card}\mathrm{Fl}(A)$, \deleted{et} $\operatorname{card\,red}(A)=\operatorname{card}(A_0)$, si $A_0$ est une catégorie réduite équivalente à $A$.

% element: text | upper third, second lemma clause | Soit C zero subset M
2°) Soit $C_0\subset M$ une sous-catégorie \inserted{\unsure{pl|de}} d'une catégorie $M$, \inserted{$\pi_0$ un cardinal infini,} et soit $\unsure{C'_0|C_0}$ la sous-catégorie str. pleine de $M$ \deleted{formée} engendrée par les limites inductives \deleted{$\varinjlim$ d'obj} dans $\deleted{\unsure{C|E}}M$ (\unsure{représentables|représentées} dans $\deleted{\unsure{C|E}}M$) des foncteurs
% element: diagram | middle, indexing functor | I to C zero
\diagram{\begin{tikzcd}[ampersand replacement=\&]
I\arrow[r,two heads]\&C_0,
\end{tikzcd}}
% element: text | middle, cardinal bounds | pour I categorie telle que
\unsure{pour|avec} $I$ catégorie telle que $\operatorname{card}\bar I\leq\pi_0$. Si $C_0$ est équivalente à une petite catégorie, il en est de même de $\deleted{\unsure{C|C_0|E}}C'_0$. Plus précisément, $\operatorname{card\,red}(C_0)\leq\pi_1$, \inserted{$\pi_1$ un cardinal infini,} alors \deleted{card} $\operatorname{card\,red}C'_0\leq\pi_1^{\pi_0}$.
% element: margin note | left lower middle, slanted | Cette majoration
% slanted
\margin{Cette majoration n'est pas prouvée, et \unsure{pro.b.blement|prob.blement|probablement} \unsure{fausse} \unsure{si on ne fait pas} des hypothèses un peu plus \unsure{strictes}}

% element: text | lower third, corollary heading and statement | Corollaire 4.6.1.2
\deleted{\unsure{3°)|3)}} \uline{Corollaire} 4.6.1.2. $M$ et $C_0$ étant comme dans 2°) ci-dessus, soit $\unsure{C_0|C}$ \unsure{le|la} plus petite sous-catégorie str. pleine de $M$ \deleted{telle que} contenant $C_0$, et stable par $\unsure{\underline{\lim}|\varinjlim}$ de cardinal $\leq\pi_0$ (celles qui sont repr. dans $M$). Si $C_0$ est équivalente à une petite catégorie, il en est de $\unsure{=|\bar m}$ de $C$. Plus précisément, si $\pi_1$ est un cardinal infini, \inserted{et $\operatorname{card\,red}C_0\leq\unsure{\pi_1|\pi}$,} alors
